\section*{Chapitre \ref{ch:iv2:vol-surface-and-vol-trading} --- La surface de volatilité et le trading de volatilité}
\markboth{\protect\omTocMark{}{Réponses~: La surface de volatilité et le trading de volatilité}}{\protect\omTocMark{}{Réponses~: La surface de volatilité et le trading de volatilité}}

\begin{solution}{iq:iv2:vol-surface-and-vol-trading:1}
C'est environ 22.4~\%.

La somme des carrés des rendements vaut 10, en unités de pourcentage au carré. La variance quotidienne vaut donc $10/5 = 2$, et la volatilité quotidienne vaut environ 1.414~\%. Multiplier par $\sqrt{252}$.

Les moyennes quotidiennes sont si petites devant la volatilité quotidienne que l'estimation de la moyenne ajoute du bruit et change peu de chose. Une moyenne de 5 points de base représente un vingt-cinquième d'une volatilité quotidienne de 1.25~\%.

\iqwrong{Multiplier par 252. Les variances sont proportionnelles au temps, et les volatilités sont proportionnelles à sa racine carrée.}
\end{solution}

\begin{solution}{iq:iv2:vol-surface-and-vol-trading:2}
Environ 1.5~\% par jour. Le 16 correspond à $\sqrt{252}$, soit environ 15.9. En comptant tous les jours calendaires, on a $\sqrt{365} \approx 19$, ce qui donne environ 1.26~\%.

\iqwrong{Prendre 1.5~\% pour le mouvement attendu. C'est un écart-type. Le mouvement absolu attendu vaut 0.8 fois ce chiffre, soit environ 1.2~\%.}
\end{solution}

\begin{solution}{iq:iv2:vol-surface-and-vol-trading:3}
Il faut environ 1.57~\$, soit un écart-type quotidien. Dans le modèle, le titre dépasse ce niveau environ 32~\% du temps.

Le gain d'une journée vaut $\tfrac12\Gamma m^2$ et le theta d'une journée vaut $\tfrac12\Gamma S^2\sigma^2\Delta t$. Ils sont égaux lorsque $m = S\sigma\sqrt{\Delta t} = 100 \times 0.25/\sqrt{252} \approx 1.57$. On perd donc environ deux jours sur trois, de petits montants, et on gagne un jour sur trois, pour des montants plus importants.

\iqwrong{S'attendre à gagner de l'argent plus d'un jour sur deux. Une position longue en gamma perd la plupart des jours, même lorsqu'elle est évaluée équitablement.}
\end{solution}

\begin{solution}{iq:iv2:vol-surface-and-vol-trading:4}
C'est environ 18.5~\%. La variance, pondérée par le temps, s'additionne d'une période indépendante à l'autre, ce que ne fait pas la volatilité.

$0.22^2 \times 0.5 = 0.0242$, et $0.25^2 \times 0.25 = 0.0156$. La variance forward vaut $(0.0242 - 0.0156)/0.25 = 0.0343$. Sa racine carrée vaut 18.5~\%.

\iqwrong{Répondre 19~\%, à partir de $2 \times 22 - 25$. Cela revient à moyenner les volatilités~; le résultat est proche ici, mais s'écarte davantage pour une courbe plus pentue.}
\end{solution}

\begin{solution}{iq:iv2:vol-surface-and-vol-trading:5}
La variance totale décroît avec le temps, ce qui est impossible. Elle impliquerait une variance forward négative. Vendre l'option d'achat à trois mois et acheter celle à six mois rapporte environ 0.34~\$.

La variance totale vaut $0.09 \times 0.25 = 0.0225$ à trois mois, et $0.04 \times 0.5 = 0.02$ à six mois. L'option d'achat à trois mois vaut environ 5.98, et celle à six mois environ 5.64. Avec des taux nuls et sans dividende, la plus longue ne peut valoir que davantage. Ainsi, à la première échéance, on peut toujours vendre l'option d'achat la plus longue pour plus que ce qu'on doit sur la plus courte.

\iqwrong{Dire que la volatilité à court terme peut être supérieure à la volatilité à long terme. Elle peut l'être, et c'est souvent le cas. Ici, elle est trop élevée, même en variance totale.}
\end{solution}

\begin{solution}{iq:iv2:vol-surface-and-vol-trading:6}
Il faut commencer par là~: les rendements des indices ne sont pas lognormaux. Les fortes baisses sont plus probables que les fortes hausses, et la volatilité augmente lors d'une baisse. Ensuite, la demande~: les investisseurs veulent des puts pour assurer leurs portefeuilles, et les dealers qui les vendent facturent ce risque. Enfin, la corrélation augmente lors d'une baisse, si bien que la volatilité des indices monte encore davantage que celle des actions.

Le modèle à volatilité plate ne comporte qu'une seule volatilité constante. Un skew traduit le fait que la distribution risque-neutre possède une queue gauche épaisse. Historiquement, les skews d'indice ne sont devenus pentus qu'après le krach de 1987.

Pour construire sa propre réponse, il faut partir de la forme de la distribution, puis exposer l'offre et la demande qui la sous-tendent, et enfin dire ce que l'on testerait. Par exemple, on peut se demander si le skew est plus pentu pour un indice que pour ses composantes, ce qui est le cas.

\iqwrong{Dire que le marché est en désaccord avec Black-Scholes parce que le modèle est faux, et s'en tenir là. Le recruteur veut savoir en quoi il est faux, et qui paie quoi.}
\end{solution}

\begin{solution}{iq:iv2:vol-surface-and-vol-trading:7}
Le notionnel de variance vaut 2\,500~\$ par point de variance. On gagne 562\,500~\$ à 25, et on perd 437\,500~\$ à 15. Le gain est plus grand, car le paiement porte sur le carré de la volatilité.

Le notionnel de variance est $100\,000/(2 \times 20)$. Le payoff vaut $2\,500(\sigma^2 - 400)$. Donc $625 - 400 = 225$, et $225 - 400 = -175$. Au voisinage du strike, un point de volatilité vaut 100\,000~\$, ce qui explique le nom de notionnel vega.

\iqwrong{Traiter le payoff comme linéaire en volatilité. C'est un swap de volatilité. Un swap de variance est convexe, et le vendeur le ressent lors d'un krach.}
\end{solution}

\begin{solution}{iq:iv2:vol-surface-and-vol-trading:8}
Il s'agit d'environ 3.2 points, et il faut 200 jours pour 1 point.

La règle est $\sigma/\sqrt{2n}$. Donc $20/\sqrt{40} = 3.16$. Pour 1 point, $2n = 400$. Sur 20 jours, une vraie volatilité de 20~\% peut ressembler à une valeur comprise entre environ 14 et 26.

\iqwrong{Comparer une volatilité réalisée sur 20 jours à l'implicite et prendre un écart de 2 points pour un signal. Un écart de cette taille reste dans le bruit.}
\end{solution}

\begin{solution}{iq:iv2:vol-surface-and-vol-trading:9}
C'est 36~\%. Pour un grand indice à volatilités égales et de corrélation deux à deux $\rho$, $\sigma_I^2 \approx \rho\sigma^2$. Donc $\rho = (18/30)^2$.

Une hausse signifie que le marché des options fait payer plus cher le mouvement conjoint des actions, ou que les options sur indice sont devenues plus chères que les options sur actions. C'est souvent un signe de peur, mais il faut vérifier que la demande de protection sur indice n'a pas simplement augmenté.

\iqwrong{Utiliser $18/30 = 60~\%$. La corrélation multiplie la variance, et il faut en prendre le carré.}
\end{solution}

\begin{solution}{iq:iv2:vol-surface-and-vol-trading:10}
À 100, la volatilité implicite vaut environ 19.98~\%. À 80, elle vaut environ 22.91~\%, et à 130, environ 23.60~\%. L'incertitude sur la volatilité crée donc un smile. À la monnaie, on obtient à peu près la moyenne, car un prix à la monnaie est presque linéaire en volatilité.

La moyenne des deux volatilités vaut 20~\%, et la moyenne quadratique vaut 22.36~\%. Pour un strike lointain, le prix est convexe en volatilité~: le mélange vaut donc plus que la moyenne, et sa volatilité implicite est plus élevée. C'est ce qui produit un smile avec une volatilité de la volatilité et sans skew~: un mélange symétrique donne une forme symétrique en log-strike, tant qu'il n'y a pas de corrélation avec le titre.

\iqwrong{Penser qu'un smile prouve que le modèle comporte des sauts. Une volatilité aléatoire seule peut y suffire, si bien qu'un smile ne dit pas quelle explication est la bonne.}
\end{solution}

\begin{solution}{iq:iv2:vol-surface-and-vol-trading:11}
À partir de ces prix arrondis, on obtient environ 0.0196 par dollar, proche de la valeur vraie 0.01985. Une réponse négative serait un arbitrage.

Le butterfly donne $(C_{95} - 2C_{100} + C_{105})/5^2$. La densité est la dérivée seconde du call par rapport au strike\index{Breeden-Litzenberger}. La densité lognormale en $K$ vaut $\varphi(d_2)/(K\sigma\sqrt T)$, soit ici $\varphi(-0.1)/20$.

L'arrondi compte ici. Avec des prix arrondis au cent, l'erreur sur la somme peut atteindre 2 cents, soit 0.0008 dans la densité, ce qui représente 4~\% de la réponse. Avec des prix non arrondis, le butterfly donne 0.0198.

\iqwrong{Utiliser des strikes trop rapprochés sur des cotations arrondies. Les butterflies plus étroits captent davantage d'arrondi, et le bruit peut faire paraître la densité négative alors qu'elle ne l'est pas.}
\end{solution}

\begin{solution}{iq:iv2:vol-surface-and-vol-trading:12}
Sticky strike~: la volatilité à la monnaie vaut 19.5~\%, et le strike 100 reste à 20~\%. Sticky moneyness~: la volatilité à la monnaie reste à 20~\%, et le strike 100 passe à 20.5~\%. Le profit est d'environ 20 cents plus élevé sous sticky moneyness.

Sous sticky strike, chaque strike conserve sa propre volatilité, si bien que la nouvelle volatilité à la monnaie est celle de l'ancien strike 105~: $20 - 0.5$. Sous sticky moneyness, tout le smile suit le titre, de sorte que le strike 100 se trouve désormais hors de la monnaie de 5~\%~: $20 + 0.5$. L'écart vaut un demi-point, à 0.39 par point, soit environ 0.2.

\iqwrong{Dire que l'un de ces deux modèles est vrai. Ce sont deux extrêmes de ce que fait le marché, et le plus proche de la réalité dépend de l'ampleur et de la vitesse du mouvement.}
\end{solution}

\begin{solution}{iq:iv2:vol-surface-and-vol-trading:13}
Le delta vaut 0.54 sous sticky strike, environ 0.62 sous sticky moneyness et environ 0.46 sous le modèle de volatilité locale. Aucun n'est valable dans tous les marchés. Il faut donc mesurer empiriquement le mouvement de la surface de volatilité en fonction du titre, et couvrir en conséquence.

La règle générale est $\Delta = \Delta_{\text{BS}} + \text{vega}\,\partial\sigma/\partial S$. Sous sticky strike, le dernier terme est nul. Sous sticky moneyness, $\sigma$ dépend de $K/S$, si bien qu'une hausse de $S$ relève la volatilité à strike fixé de 0.002 par dollar. Le terme supplémentaire vaut $39.7 \times 0.002 = 0.079$. La volatilité locale agit dans le sens inverse~: $0.54 - 0.079$.

\iqwrong{Citer un seul delta et croire qu'il ne repose sur aucun modèle. Un delta contient une hypothèse sur la façon dont le smile se déplace.}
\end{solution}

\begin{solution}{iq:iv2:vol-surface-and-vol-trading:14}
Un point sur les 1.5 provient du glissement le long du smile, et un demi-point du déplacement du smile lui-même vers le haut. L'option au strike à la monnaie initial se trouve 2~\% au-dessus du nouveau niveau à la monnaie, si bien que sa volatilité est passée de l'ancien niveau à la monnaie à environ un demi-point au-dessus de celui-ci.

Le nouveau strike à la monnaie est celui qui se trouvait 2~\% en dessous, si bien que le glissement donne $2 \times 0.5 = 1$. Les 0.5 restants correspondent à un déplacement parallèle. L'ancien strike ne bouge que de ce déplacement parallèle, puisqu'il reste à sa place le long de la courbe initiale.

\iqwrong{Dire que la volatilité implicite de l'option a monté de 1.5 point. Elle l'a fait au nouveau strike à la monnaie, et le strike de l'option n'est pas celui-ci.}
\end{solution}

\begin{solution}{iq:iv2:vol-surface-and-vol-trading:15}
On gagne environ 363\,600~\$. Cela représente environ 3.6 points de vega. La volatilité implicite peut rester au-dessus de la volatilité réalisée, car les investisseurs paient pour se protéger d'une hausse de la volatilité lors d'une chute, et les vendeurs veulent être rémunérés pour ce risque. Le piège est que le vendeur perd de fortes sommes au pire moment.

Le payoff est $100\,000/(2 \times 22) \times (484 - 324) = 2\,273 \times 160$. Le strike était supérieur de quatre points à la volatilité réalisée, mais la convexité fait qu'il vaut un peu moins que quatre.

Au vu des données, les options sur indices couvertes en delta ont perdu de l'argent en moyenne, résultat connu sous le nom de prime de risque de volatilité négative. Un jour de février 2018, le S\&P 500 a chuté de 4~\% et son indice de volatilité a grimpé de 20 points. Cet écart n'est donc pas un gain sans risque.

\iqwrong{Traiter l'écart moyen comme un gain sans risque. Le fait qu'il paie dépend de l'ampleur des mouvements et de la queue de distribution.}
\end{solution}

\begin{solution}{iq:iv2:vol-surface-and-vol-trading:16}
L'estimateur est celui de Parkinson\index{estimateur de Parkinson}~: le carré de l'amplitude logarithmique, multiplié par $1/(4\ln 2) \approx 0.36$. Pour un mouvement brownien sans dérive, $E[\text{amplitude}^2] = 4\ln 2\,\sigma^2$. La variance de cet estimateur vaut environ un cinquième de celle d'un rendement au carré, soit environ cinq jours de données close-to-close pour une seule journée.

Les biais~: les prix réels ne sont observés qu'à des instants discrets, et le vrai plus haut n'est jamais observé, ce qui abaisse l'estimation. L'écart overnight est totalement laissé de côté. Le rebond bid-ask gonfle l'amplitude.

Il faut donc utiliser une forme qui ajoute un terme overnight, et comparer avec une volatilité réalisée calculée à partir de rendements intrajournaliers.

\iqwrong{Utiliser sans le facteur correctif. Le carré de l'amplitude a une espérance d'environ 2.77 fois la variance.}
\end{solution}

\begin{solution}{iq:iv2:vol-surface-and-vol-trading:17}
L'écran affiche environ 17.4~\%. Ce chiffre paraît bas, car l'écran compte trois jours de temps, alors que le titre n'est ouvert que pendant un seul. Le lundi matin, le même écran affiche environ 30~\%.

En égalant les variances, $\sigma_c^2 \times 3/365 = 0.25^2/252$. Donc $\sigma_c = 25~\% \times \sqrt{365/756} \approx 17.4~\%$.

Le lundi, un jour ouvré vaut un jour calendaire. Alors $25~\% \times \sqrt{365/252} \approx 30~\%$. Le très court terme d'une courbe en temps calendaire présente donc une dent de scie hebdomadaire, et ce motif n'a pas de signification en soi. Un trader qui ne voit pas cela trouvera de fausses opportunités.

\iqwrong{Acheter la volatilité du vendredi parce qu'elle paraît bon marché, ou vendre celle du lundi parce qu'elle paraît chère. Utiliser une horloge qui donne aux jours ouvrés et aux jours d'événement leur poids réel.}
\end{solution}

\begin{solution}{iq:iv2:vol-surface-and-vol-trading:18}
La volatilité de base vaut environ 33.2~\%, et le mouvement implicite de l'événement environ 7.0~\%, si bien que le mouvement absolu attendu vaut environ 5.6~\%. On suppose que la volatilité de base est constante, et que l'événement correspond à un seul rendement indépendant.

La variance totale vaut $0.60^2 \times 5/252 = 0.00714$ pour l'option courte, et $0.40^2 \times 25/252 = 0.01587$ pour la longue. Leur différence est la variance de base sur 20 jours. Donc la variance de base vaut $0.00873 \times 252/20 = 0.110$, dont la racine est 33.2~\%. La variance d'événement vaut $0.00714 - 0.110 \times 5/252 = 0.00496$, dont la racine est 7.04~\%. Le mouvement absolu porte le facteur $\sqrt{2/\pi} = 0.8$.

\iqwrong{Prendre les 60~\% pour la volatilité de l'événement. Une grande partie de la variance de l'option courte vient de l'événement, mais les jours de base y figurent aussi.}
\end{solution}

\begin{solution}{iq:iv2:vol-surface-and-vol-trading:19}
La volatilité implicite vaut environ 32.6~\%. Après l'événement, le straddle vaut environ 5.62 si le titre reste inchangé, 6.20 à 103, 7.66 à 106 et 10.60 à 110. On atteint le seuil de rentabilité pour un mouvement d'environ 5.7~\$.

La variance totale s'additionne~: $0.25^2 + 0.06^2/(30/365)$, dont la racine vaut 32.6~\%. Le straddle coûte 7.46 à ce niveau. Après une journée, sa volatilité retombe à 25~\%, ce qui lui retire environ 1.83~\$ même sans mouvement~: c'est le vol crush.

Un straddle est donc un pari sur un mouvement supérieur au mouvement absolu attendu par le marché, soit 4.8~\%. Il en faut près de 6~\%. En raison de la valeur restante de l'option après l'événement, le seuil de rentabilité est inférieur au prix du straddle, soit 7.46.

\iqwrong{Comparer le mouvement attendu au seul prix du straddle. Le straddle garde une valeur après l'événement, il faut donc valoriser ce qui reste.}
\end{solution}

\begin{solution}{iq:iv2:vol-surface-and-vol-trading:20}
Le strike de variance juste vaut $\sqrt{0.04 + 0.0025} \approx 20.6$. Le strike de volatilité juste vaut 20. Un swap de volatilité paie moins lors d'un gros mouvement, car il porte sur la volatilité, qui est linéaire par rapport à sa propre valeur, et non sur son carré.

Le strike de variance est la racine carrée de la moyenne de la variance~: $E[\sigma^2] = 20^2 + 5^2 = 425$. Le swap de volatilité n'est que la moyenne de la volatilité.

Ce swap est difficile à couvrir, car la racine carrée est concave par rapport à la variance. Un ensemble statique d'options réplique la variance, et non la volatilité. Un swap de volatilité exige donc un modèle de volatilité de la volatilité et une gestion dynamique de la variance. L'écart entre les deux strikes, ici environ 0.6, est l'ajustement de convexité\index{ajustement de convexité}.

\iqwrong{Égaliser les deux strikes. Cela revient à ignorer la variance de la volatilité, or c'est précisément ce qu'un swap de variance fait payer.}
\end{solution}

\begin{solution}{iq:iv2:vol-surface-and-vol-trading:21}
La volatilité de l'indice vaut environ 16.8~\%. Si la corrélation se réalise à 0.2, la volatilité de l'indice tombe à environ 13.9~\%, et l'on gagne 88 points de variance sur l'indice sans changement sur les actions. Avec une volatilité des actions de 33~\%, l'indice se situe à environ 15.3~\%, et l'on gagne aussi sur les actions. C'est ce qu'on appelle le dispersion trading\index{dispersion trading}. On est vendeur de corrélation, et cela fait mal lors d'un krach, quand la corrélation bondit.

La formule est $\sigma_I^2 = \sigma^2[1/n + (1 - 1/n)\rho]$. Avec $n = 50$, $\sigma = 30~\%$ et $\rho = 0.3$, elle vaut 0.02826. Pour $\rho = 0.2$, elle vaut 0.01944, soit 88 points de moins. La même formule à 33~\% donne 15.3~\%.

On encaisse la corrélation implicite moins la réalisée, ce qui constitue une prime de risque, et l'on la paie en période de crise. Cette stratégie est aussi coûteuse à mettre en œuvre, car il faut 50 positions sur actions, et chacune doit être couverte.

\iqwrong{Qualifier la stratégie de pur trade de volatilité. C'est surtout une position vendeuse de corrélation. Même avec des pondérations choisies pour être neutres en vega, on reste exposé aux niveaux de volatilité.}
\end{solution}

\begin{solution}{iq:iv2:vol-surface-and-vol-trading:22}
Avec le skew, il coûte environ 2.69~\$. À volatilité plate de 20~\%, il coûterait environ 4.38~\$. Il est moins cher, car le put que l'on vend a une volatilité plus élevée, donc un prix plus élevé~: environ 5.27, contre 3.59.

Le put de strike 100 à 20~\% coûte 7.97. Donc $7.97 - 5.27 = 2.69$ et $7.97 - 3.59 = 4.38$.

Celui qui achète un put à la baisse paie le skew, et celui qui vend un put plus éloigné en récupère une grande partie. C'est pourquoi un investisseur qui veut une protection bon marché achète un put spread, et renonce à la queue de distribution.

\iqwrong{Utiliser une seule volatilité pour tous les strikes d'un spread. Le spread est la partie du smile où le skew compte le plus.}
\end{solution}

\begin{solution}{iq:iv2:vol-surface-and-vol-trading:23}
Le modèle ne vérifie pas la condition. Elle s'écrit $2\kappa\theta \geq \xi^2$, et ici $0.16 < 0.25$. La variance peut donc atteindre zéro, et une simulation simple la fera passer en négatif.

Si l'on simule par Euler, il faut utiliser la troncature complète~: on applique $\max(v, 0)$ à la dérive et à la diffusion. Ou bien utiliser le schéma quadratique-exponentiel, qui reproduit mieux les moments exacts. La plupart des modèles calibrés ne vérifient pas la condition, on ne peut donc pas la passer sous silence.

\iqwrong{Imposer la condition dans la calibration. Cela modifie le smile que l'on peut ajuster, alors qu'elle n'est pas nécessaire à l'existence du prix.}
\end{solution}

\begin{solution}{iq:iv2:vol-surface-and-vol-trading:24}
Interpoler linéairement la variance totale, sur l'horloge à laquelle on fait confiance. On obtient environ 21.5~\%.

Les variances totales valent $0.04/12 = 0.00333$ et $0.0484 \times 0.25 = 0.0121$. À mi-chemin dans le temps, on obtient 0.00772, pour deux mois. Multiplier par 6, puis prendre la racine.

L'interpolation linéaire en volatilité donne 21~\%, et elle peut faire décroître la variance totale sur une courbe pentue. L'interpolation linéaire en variance totale ne le peut pas, tant que les données d'entrée sont elles-mêmes croissantes. Il faut donc vérifier que la variance totale, à chaque strike et à moneyness fixée, croît avec l'échéance. Avec des événements, on pondère l'horloge, et l'on ne traverse pas ceux-ci en ligne droite.

\iqwrong{Traiter cela comme une question portant sur un seul chiffre. Sur toute la surface, la même règle doit valoir à chaque moneyness, sinon l'arbitrage calendaire réapparaît sur les ailes.}
\end{solution}

\begin{solution}{iq:iv2:vol-surface-and-vol-trading:25}
Il vaut environ 386\,250~\$. La valorisation est simple, puisque la variance s'additionne dans le temps. Son vega est tombé à environ la moitié de ce qu'il était, car il ne reste que la moitié du temps.

La variance totale espérée vaut $\tfrac12 \times 625 + \tfrac12 \times 484 = 554.5$. Donc la valeur vaut $2\,500 \times (554.5 - 400) = 386\,250$, à taux nuls.

Le vega, par point de variance restante, vaut $2\,500 \times \tfrac12$, par point de variance. À une volatilité de 22, un point de volatilité vaut $2 \times 22$ points de variance. Donc le vega vaut environ 55\,000~\$, contre 100\,000~\$ au départ. Il est linéaire par rapport au temps restant, si bien qu'un swap de variance ne réserve aucune surprise sur le theta, et c'est une exposition nette.

\iqwrong{Ne valoriser le swap que sur la volatilité implicite courante. Une part majeure de la valeur d'un swap de variance déjà entamé provient de ce qui est déjà réalisé.}
\end{solution}

\begin{solution}{iq:iv2:vol-surface-and-vol-trading:26}
Un strip pondéré par $1/K^2$ a le même dollar gamma quel que soit le cours de l'action. Son gain par unité de variance ne dépend donc pas de l'endroit où va l'action. Avec des poids égaux, le dollar gamma croît avec le cours de l'action. Le strike équitable vaut ici 20~\%, puisqu'un smile plat n'a rien d'autre à offrir.

Une option couverte en delta rapporte $\tfrac12S^2\Gamma(\sigma_r^2 - \sigma_i^2)\,dt$. Pour une exposition propre à la variance, $S^2\Gamma$ doit donc être constant. Le gamma de chaque option dépend de $S/K$ d'une façon qui fait fonctionner cela exactement lorsque $f^{\prime\prime}(K) \propto 1/K^2$.

Dans un contrôle numérique pour des cours de 80, 100 et 120, avec des strikes allant de 10 à 1\,000, le dollar gamma du strip reste identique à 0.2~\% près. Celui à poids égaux croît avec le carré du cours de l'action.

Le strip est le développement au second ordre de $-2\ln S$. C'est pourquoi il s'agit du contrat log\index{contrat log} déguisé, et c'est pourquoi le strike équitable vaut $-(2/T)E[\ln(S_T/S_0)]$.

\iqwrong{Couvrir un swap de variance par un straddle à la monnaie. Il n'offre la bonne exposition qu'à un seul cours de l'action, et cette exposition s'estompe dès que l'action bouge.}
\end{solution}

\begin{solution}{iq:iv2:vol-surface-and-vol-trading:27}
Un tel smile ne peut pas être juste, car un call gagnerait de la valeur à des strikes plus élevés. La pente maximale admise est de 2 en variance totale par unité de log-strike. Pour une action à 100, à un strike de 400, la volatilité implicite la plus élevée vaut environ 167~\%.

La formule des moments de Roger Lee\index{formule des moments de Lee} indique que la variance totale peut croître au plus comme $2|k|$ pour un $|k|$ grand, où $k = \ln(K/S)$. Au-delà, le prix de Black-Scholes avec cette variance totale ne tend pas vers zéro quand $K$ croît. Le prix d'un call doit décroître avec le strike. Et pour un $K$ grand, il doit tendre vers zéro.

Dans un contrôle, une variance totale de $3|k|$ fait croître la valeur du call avec le strike jusqu'à très loin, passant de 42 pour un strike de 300 à 82 pour $10^8$. Avec $1.5|k|$, elle décroît loin en strike, de 23 pour 1\,000 à 17 pour $10^8$, comme il se doit.

Pour l'application numérique~: $k = \ln 4 = 1.386$, donc la variance totale vaut au plus 2.77, et sa racine vaut 1.665 pour un an.

\iqwrong{Extrapoler un polynôme dans les ailes. Le faire passer par les cotations rend absurdes les prix des options lointaines, et l'erreur se répercute alors sur le strike du swap de variance, qui intègre ces options.}
\end{solution}

\begin{solution}{iq:iv2:vol-surface-and-vol-trading:28}
La formule est $\sigma_{\text{loc}}^2(K, T) = 2\partial_TC/(K^2\partial_{KK}C)$. Pour une surface plate, on retrouve 20~\%. Pour le mélange à un strike de 100, on obtient environ 17.3~\%. Elle diffère autant parce qu'il s'agit d'une moyenne conditionnelle, alors que la volatilité implicite est une moyenne sur les trajectoires.

Pour le cas plat, $\partial_TC = S\varphi(d_1)\sigma/(2\sqrt T)$, et $K^2\partial_{KK}C = K\varphi(d_2)/(\sigma\sqrt T)$. En divisant l'un par l'autre, on obtient $\sigma^2$.

La formule de Dupire\index{formule de Dupire} égale la variance instantanée espérée, sachant que l'action se trouve à ce niveau à cet instant. Dans le mélange, l'action à 100 après un an se trouve beaucoup plus probablement dans l'état à 10~\%, puisque sa densité y est plus haute. Les poids valent donc environ 75~\% et 25~\%, et $\sqrt{0.75 \times 0.01 + 0.25 \times 0.09} \approx 17.3\%$. La volatilité implicite à la monnaie vaut à peu près la moyenne des deux états.

\iqwrong{Lire la volatilité locale comme une prévision de la volatilité réalisée future. Elle est ce qui reproduit les prix d'options du jour. Ce n'est pas une estimation de ce que fera l'action.}
\end{solution}

\begin{solution}{iq:iv2:vol-surface-and-vol-trading:29}
Les volatilités implicites valent environ 22.05~\%, 19.98~\% et 18.13~\%, par simulation. Cela donne une pente d'environ $-0.1$ par unité de log-strike, soit la moitié de la pente locale. La règle veut que la pente implicite soit la moitié de la pente locale pour une courte échéance.

La volatilité implicite à un strike est proche d'une moyenne de la volatilité locale le long de la trajectoire la plus probable du spot jusqu'à ce strike. Si cette trajectoire suit une droite en log-prix, la moyenne se situe à mi-chemin. La pente implicite vaut donc la moitié de la pente locale.

Le test est une simulation sous le modèle. Sa pente est à 0.01 près de celle que donne la règle, y compris à un an.

\iqwrong{Ajuster une volatilité locale dont la pente est égale à celle du skew observé. L'ajustement est alors faux d'un facteur deux, et les prix vanille du modèle sont décalés.}
\end{solution}

\begin{solution}{iq:iv2:vol-surface-and-vol-trading:30}
La volatilité à la monnaie a baissé d'environ 0.97 point, soit environ deux fois le déplacement propre du skew, de 0.49, sur ce mouvement. Et la volatilité du strike 100 a baissé d'un demi-point, si bien que le smile s'est déplacé vers la gauche tandis que le cours de l'action montait. Les marchés ont plutôt tendance à faire l'inverse, en déplaçant le smile avec l'action.

Le delta du modèle a donc la mauvaise correction. Il retranche un terme de vega là où le marché exige qu'on l'ajoute. Ses couvertures sont donc moins bonnes que celles de Black-Scholes dans de nombreux tests, ce qui constitue la critique bien connue de Hagan et de ses coauteurs.

La volatilité locale reste utile. Elle réévalue exactement les options vanille, et elle donne aux produits de type barrière une valeur cohérente avec chaque prix vanille. Mais la dynamique est quelque chose qu'il faut ajouter soi-même.

\iqwrong{Croire qu'un modèle qui ajuste toutes les options doit bien les couvrir. Un ajustement est un constat statique, alors qu'une couverture porte sur la façon dont cet ajustement change.}
\end{solution}

\begin{solution}{iq:iv2:vol-surface-and-vol-trading:31}
La structure par terme vaut environ 22.5~\%, 26.2~\% et 28.9~\%, en tendant vers $\sqrt{0.09} = 30~\%$. La corrélation fixe la pente, et la vol of vol fixe la courbure. La volatilité à la monnaie est inférieure à 20~\%, car les prix à la monnaie sont presque linéaires en volatilité, et la moyenne d'une volatilité est inférieure à la racine de sa variance moyenne.

La variance moyenne espérée sur $T$ vaut $\theta + (v_0 - \theta)(1 - e^{-\kappa T})/(\kappa T)$. Donc, à un an, $0.09 - 0.05 \times 0.432 = 0.0684$, et la racine carrée vaut 26.2~\%. Une simulation confirme les 26.2~\%.

Pour le smile~: avec une corrélation de $-0.7$, une action qui baisse fait monter la volatilité, ce qui donne une queue gauche plus épaisse, et c'est cela le skew. Avec une corrélation nulle, le smile est symétrique en log-strike. Sa hauteur est donnée par la vol of vol.

\iqwrong{Croire que la vitesse de retour à la moyenne fixe le skew. Elle règle la vitesse à laquelle le skew s'estompe avec l'échéance.}
\end{solution}

\begin{solution}{iq:iv2:vol-surface-and-vol-trading:32}
À la monnaie, elle vaut environ 20.11~\%. À 80, elle vaut environ 21.93~\%, et à 120, environ 19.44~\%. Une simulation donne 20.14, 21.87 et 19.44. Son avantage est que le smile se déplace dans le bon sens avec l'action. Sa faiblesse est que son développement peut donner des densités négatives pour les bas strikes et les longues échéances.

La formule à la monnaie, avec $\beta = 1$, s'écrit $\alpha[1 + (\rho\nu\alpha/4 + (2 - 3\rho^2)\nu^2/24)T]$. Donc $0.2 \times (1 - 0.006 + 0.01153) = 0.2011$, soit 20.11~\%. Hors de la monnaie, la formule utilise $z = (\nu/\alpha)\ln(F/K)$.

C'est une formule et non la solution d'un modèle~: elle peut donc se tromper pour les bas strikes et les longues échéances. Cela compte surtout en taux, où elle sert d'outil de cotation du marché. Il faut la vérifier par simulation, ou par comparaison avec une densité calculée à partir des prix donnés par la formule elle-même.

\iqwrong{Utiliser le modèle SABR pour valoriser des produits dépendant du chemin. C'est un bon interpolateur, et il donne un delta. Il ne s'ajuste pas à la structure par terme dans son ensemble.}
\end{solution}

\begin{solution}{iq:iv2:vol-surface-and-vol-trading:33}
À un mois, les volatilités implicites à 90, 100 et 110 valent environ 20.8~\%, 15.6~\% et 15.2~\%. À un an, elles valent environ 16.7~\%, 16.3~\% et 16.0~\%. Le skew est si pentu à courte échéance parce qu'une diffusion bouge à peine en un mois, de sorte que seul le saut peut atteindre 90. Le skew d'une diffusion pure s'aplatit aux courtes échéances, sauf si sa volatilité est extrême près de ces strikes.

Le call s'écrit $\sum_jp_jC_{\text{BS}}(F_j, K, \sigma, T)$, où les $p_j$ suivent une loi de Poisson de moyenne $\lambda T$, et $F_j = 100e^{-\lambda JT}(1 + J)^j$ avec $J = -0.2$. Le facteur $e^{-\lambda JT}$ rétablit la perte espérée due aux sauts, afin que l'action soit une martingale.

Ainsi, à l'extrémité courte de la courbe, le skew tient surtout au saut, et son niveau dépend de la probabilité et de la taille de ce saut. Sur un an, le saut se fond dans le mélange, et le skew s'aplatit vite. La volatilité locale peut reproduire cette surface aujourd'hui, mais seulement avec des niveaux de volatilité extrêmes près de l'échéance.

\iqwrong{Chercher à reproduire un skew à un mois très pentu avec la seule volatilité stochastique diffusive. Il faudrait une vol of vol extrême ou un retour à la moyenne très rapide, et le modèle se brise alors aux échéances plus longues.}
\end{solution}

\begin{solution}{iq:iv2:vol-surface-and-vol-trading:34}
Un saut à la baisse de 20~\% a un rendement log au carré de 0.0498, et le strip donne 0.0463, de sorte que l'on perd 0.0035 par unité de notionnel de variance. Un saut à la hausse de 20~\% a un rendement log au carré de 0.0332, et le strip donne 0.0354, de sorte que l'on gagne 0.0021.

Le strip verse $2(r - \ln(1 + r))$ sur un rendement simple $r$, puisqu'il s'agit du contrat log diminué d'une couverture en delta. Cela ne coïncide avec $r^2$ que pour de petits mouvements. Pour $r = -0.2$, $2(-0.2 + 0.2231) = 0.0463$. Pour $r = +0.2$, il vaut $2(0.2 - 0.1823) = 0.0354$.

Vendre un swap couvert par un strip revient donc à vendre un saut à la baisse. Ainsi, le strike équitable d'un swap de variance se situe au-dessus de ce que suggère une réplication continue, avec une distribution à skew négatif. C'est aussi pourquoi les swaps sur titres individuels sont souvent plafonnés~: la perte du vendeur en cas de défaut est énorme.

\iqwrong{Considérer la réplication d'un swap de variance comme exacte. Elle ne l'est que pour une trajectoire continue. En présence d'un saut, apparaît le terme $r^3$, et c'est son signe qui fait le dégât.}
\end{solution}

\begin{solution}{iq:iv2:vol-surface-and-vol-trading:35}
Le strike équitable issu du strip vaut environ 20.54~\%, soit 0.0422 exprimé en variance. La règle donne $0.04(1 + 0.03) = 0.0412$, soit 20.30~\%. Le strike est au-dessus du niveau à la monnaie parce que le strip inclut les puts du côté baissier, et qu'avec un skew ceux-ci sont chers.

Le strip vaut $(2/T)\int(\text{option hors de la monnaie})/K^2\,dK$, et il couvre tous les strikes. Ses poids sont plus lourds du côté baissier, et c'est dans l'aile des puts que le smile est le plus élevé. La variance se situe donc au-dessus de la monnaie, quel que soit le sens du skew, et le $b^2$ de la règle le montre.

La règle est grossière~: c'est un développement au second ordre avec un skew linéaire, et elle néglige la courbure. Elle est donc inférieure de 0.24 point ici. C'est un contrôle utile, et non un prix.

\iqwrong{Prendre le niveau à la monnaie comme strike de variance. Sur un marché d'indices, il peut se situer plusieurs points en dessous.}
\end{solution}

\begin{solution}{iq:iv2:vol-surface-and-vol-trading:36}
L'indice a une volatilité d'environ 13.4~\% dans l'état normal et d'environ 26.8~\% dans l'autre. Avec le même forward, le smile vaut environ 17.4, 16.1 et 17.0, pour des niveaux de 80, 100 et 120. C'est un smile, et presque aucun skew. Pour obtenir un skew, il faut faire en sorte que l'état à forte corrélation soit celui dans lequel les cours ont baissé.

Pour un grand indice, $\sigma_I^2 = \rho\sigma^2$. Les poids du mélange valent 0.8 et 0.2. Une vérification avec un forward de 85 dans l'état à forte corrélation et de 103.75 dans l'autre, ce qui conserve la moyenne à 100, donne un skew baissier marqué.

Les skews des indices sont plus prononcés parce qu'une baisse de l'indice va souvent de pair avec une corrélation plus forte, en plus d'une volatilité plus élevée des titres. Un titre isolé a ses propres nouvelles, bonnes ou mauvaises. Selon un résultat documenté, les distributions risque-neutre d'un titre isolé sont nettement moins asymétriques vers la gauche que celles de l'indice.

\iqwrong{Affirmer que la forte corrélation cause à elle seule le skew. Elle ajoute de la volatilité dans un état. C'est le lien entre cet état et la baisse des cours qui produit le skew.}
\end{solution}

\begin{solution}{iq:iv2:vol-surface-and-vol-trading:37}
La moyenne vaut environ 0.93, pour une prime de 4.61, et l'écart type environ 0.75. On perd environ 9~\% du temps, et le 1\up{er} centile vaut environ $-0.68$. Le ratio de Sharpe par transaction est donc d'environ 1.2, ce qui est un cas irréaliste de bonne performance.

La moyenne est la différence entre les deux prix, à 20~\% et à 16~\%, soit 0.92. La voie par le vega donne le même résultat~: $4 \times 0.23$.

Votre modèle ne comporte pas de sauts, et sa volatilité est connue. En réalité, la volatilité réalisée est incertaine, et un mouvement de grande ampleur peut effacer plusieurs mois de gains. On dimensionne donc la position à partir d'un scénario de stress~: par exemple, un écart de 20~\% sur le straddle coûte près de 15.4~\$ ici, soit environ 17 fois le profit moyen. Il reste ensuite à décider quelle part de cette perte on peut supporter.

\iqwrong{Dimensionner la position avec l'écart type du modèle. Une stratégie short gamma paraît la meilleure dans les statistiques justement pendant la période qui précède sa défaillance.}
\end{solution}

\begin{solution}{iq:iv2:vol-surface-and-vol-trading:38}
Par Itô, $\tfrac12\int\sigma^2\,dt = \int dS/S - \ln(S_T/S_0)$. La variance équitable vaut donc $-(2/T)E[\ln(S_T/S_0)]$, et le développement de Carr-Madan du logarithme donne
\[
K_{\text{var}} = \frac2T\Big[\int_0^{S_0}\frac{P(K)}{K^2}\,dK + \int_{S_0}^\infty\frac{C(K)}{K^2}\,dK\Big].
\]
La valeur exacte vaut 0.04218, soit 20.54~\%. Avec des strikes de 50 à 200, un pas de 1~\$ donne 0.04205, un pas de 5~\$ donne 0.04250 et un pas de 10~\$ donne 0.04378. De 70 à 140 par pas de 5~\$, on obtient 0.04096.

Les deux erreurs ne vont pas dans le même sens. Un pas trop grossier surestime, puisque l'intégrande est convexe, comme avec la règle des trapèzes sur un point anguleux en $S_0$. C'est le cas des valeurs 0.04250 et 0.04378. Un intervalle tronqué sous-estime, car il laisse de côté des options, et c'est le cas de 0.04096 et d'une partie de 0.04205.

En situation de stress, la troncature pèse davantage. Quand un titre baisse, une part plus grande de la variance se trouve sur des strikes situés hors de l'intervalle initial. C'est pourquoi le gamma d'un portefeuille de réplication s'effondre au-delà de son dernier strike. L'écart entre un strip coté et un strip réel est ce que découvre un fonds qui croit avoir couvert son risque sans l'avoir fait.

\iqwrong{Ne chercher que des strikes plus fins. Le coût des ailes croît avec la volatilité, alors que l'erreur due au pas est facile à réduire. Un intervalle plus large est là où se trouve l'essentiel de l'enjeu.}
\end{solution}

\begin{solution}{iq:iv2:vol-surface-and-vol-trading:39}
On conditionne sur la trajectoire de la volatilité. La action suit alors une loi lognormale de variance $\bar v = \frac1T\int v\,dt$, si bien que son prix d'option d'achat vaut $C_{\text{BS}}(\sqrt{\bar v})$. La moyenne de ce prix sur les trajectoires est le prix recherché. C'est la formule de mélange de Hull-White\index{Hull-White mixing}. Dans un test, les deux méthodes donnent le même smile à près du bruit près, et l'erreur type de l'estimation par mélange est environ 5 fois plus petite à la monnaie.

Pour le résultat du mélange, il suffit de simuler la variance et d'enregistrer $\bar v$. Une simulation directe a aussi besoin du bruit du titre.

La seconde méthode est donc une estimation de Monte Carlo conditionnelle~: elle élimine exactement le bruit du titre. C'est pourquoi elle réduit la variance, sans coût supplémentaire.

Avec la corrélation, conditionner sur $v$ modifie aussi l'espérance du titre, puisque le bruit du titre est lié à celui de la variance. Le prix conditionnel du titre suit alors une loi lognormale avec un forward différent~: on multiplie $S$ par $\exp(\rho\int\sqrt v\,dW - \tfrac12\rho^2\int v\,dt)$ et on utilise la variance restante $(1 - \rho^2)\bar v$. La forme de Romano-Touzi de la même idée traite ce cas.

\iqwrong{Utiliser un mélange avec la mauvaise condition. Si la corrélation n'est pas nulle, la formule simple ne donne aucun skew, ce qui est faux.}
\end{solution}

\begin{solution}{iq:iv2:vol-surface-and-vol-trading:40}
La variance locale vaut $\sigma^2_{\text{loc}}(x, t) = w_1\sigma_1^2 + w_2\sigma_2^2$, où chaque $w_i$ est proportionnel à la densité de l'état $i$ au logarithme du cours $x$ à l'instant $t$. Une simulation quotidienne naïve échoue~: pour un strike de 70, elle donne une volatilité implicite d'environ 23.2~\%, contre 24.97~\%. Une équation aux dérivées partielles sur une grille temporelle dont les pas se resserrent près de zéro le résout à 0.02 point près.

La formule provient du théorème de Gyöngy\index{Gyöngy's theorem}~: la variance locale vaut $E[v \mid S_t = S]$. Le poids d'un état est sa part de la densité. Grâce à cela, un seul modèle à un facteur aléatoire reproduit toutes les lois marginales du mélange.

Au voisinage de l'instant zéro, la variance locale est singulière. Les deux densités ont des largeurs très différentes~: pour de petits $t$, l'état à 30~\% prend le dessus très près de l'origine, et $\sigma_{\text{loc}}$ passe de 10~\% à 30~\% sur une distance très courte. Un pas quotidien saute par-dessus cette zone. La trajectoire manque donc la dispersion initiale, et les strikes éloignés sortent trop bas. Le prix moyen dérive également à partir de 100.

La solution consiste à prendre des pas petits au départ, par exemple $t_i = (i/m)^2$, et à utiliser une grille implicite en $x$. Comme le repricing en volatilité locale doit être exact, c'est un contrôle utile de toute calibration.

\iqwrong{Attribuer cette erreur à l'erreur d'échantillonnage. Avec davantage de trajectoires, le résultat reste faux de la même quantité. Seuls davantage de pas au début corrigent le problème.}
\end{solution}
